\begingroup
\newif\ifAdvHideHeading
\RenewDocumentEnvironment{advexamplebox}{m}{%
  \refstepcounter{advexample}%
  \begin{tcolorbox}[
    enhanced,unbreakable,colback=white,colframe=RuleGray,boxrule=0.5pt,arc=0mm,
    title={例~\theadvexample\quad #1},fonttitle=\bfseries,
    coltitle=Ink,colbacktitle=PaleBlue,borderline west={1.4pt}{0pt}{DeepBlue}]
}{\end{tcolorbox}}

\RenewDocumentEnvironment{advsolution}{+b}{%
  \par\smallskip\begingroup\small\noindent
  \ifAdvHideHeading\global\AdvHideHeadingfalse\else\LectureSolutionHeading\fi #1\par\endgroup\smallskip
}{}

\RenewDocumentEnvironment{advdefinition}{m +b}{%
  \refstepcounter{advdefinition}%
  \begin{tcolorbox}[
    enhanced,unbreakable,colback=white,colframe=RuleGray,boxrule=0.5pt,arc=0mm,
    title={定义~\theadvdefinition\quad #1},fonttitle=\bfseries,
    coltitle=Ink,colbacktitle=PaleGold,borderline west={1.2pt}{0pt}{Gold}]
  #2
  \end{tcolorbox}
}{}

\RenewDocumentEnvironment{advtheorem}{m +b}{%
  \refstepcounter{advtheorem}%
  \begin{tcolorbox}[
    enhanced,unbreakable,colback=white,colframe=RuleGray,boxrule=0.5pt,arc=0mm,
    title={定理~\theadvtheorem\quad #1},fonttitle=\bfseries,
    coltitle=Ink,colbacktitle=PaleBlue,borderline west={1.2pt}{0pt}{DeepBlue}]
  #2
  \end{tcolorbox}
}{}


\AdvancedUnit{2}{热学}

\section{热力学}

\subsection{热力学定律与基本关系}

热力学（thermodynamics）研究宏观体系的平衡态与状态变化。研究对象可以是易拉罐中的饮料、大气中的气团，也可以是外磁场中的超导体；选定对象称为\emph{体系}，其余部分称为外界。第一定律排除第一类永动机，第二定律排除把单一热源吸收的热量全部变成功的循环装置，第三定律则说明绝对零度不能由有限步骤达到。不同表述的逻辑起点虽有差别，但最终必须给出相同的可观测预言。

体系处于平衡态时，内部各子体系之间以及体系与外界之间均无宏观净流。体系具有内能 $U$；第 $i$ 种物质用粒子数 $N_i$ 表示，辐射等粒子数不守恒的成分则不必把粒子数列为独立变量。

\begin{advtheorem}{宏观态的完备参数组}\label{advA2:thm:macrostate}
一个平衡宏观态可由内能、宏观广义坐标和各组分粒子数完整描述：
\[
 A=A(U,\{X_\alpha\},\{N_i\}).
\]
若只保留一个体积坐标，则写成 $A=A(U,V,\{N_i\})$。
\end{advtheorem}
这里“完整”意味着参数空间与宏观状态空间之间是一一对应：不存在遗漏状态，同一组参数也不代表两个不同的宏观态。实际过程若不能视为连续平衡态序列，称为非准静态过程；全过程都经过平衡态的理想极限称为准静态过程。准静态还不自动等于可逆，只有不存在摩擦、有限温差传热等耗散时才可逆。

均匀体系中的量分为广延量和强度量。广延量随体系复制而同比放大，
\[
 A(\lambda U,\lambda V,\lambda N)=\lambda A(U,V,N),
\]
强度量则满足 $a(\lambda U,\lambda V,\lambda N)=a(U,V,N)$。内能、体积和粒子数既是广延量，又对互不相互作用的子体系具有加性；大量相同体系构成系综。孤立体系满足 $U,V,N$ 均不与外界交换。

\begin{advtheorem}{熵原理}\label{advA2:thm:entropy-principle}
平衡热力学可用一个加性态函数
\[
 S=S(U,V,N)
\]
组织。在本节所讨论的正温度体系中，熵随内能单调增加；孤立体系的可逆过程总熵不变，不可逆过程总熵增加。
\end{advtheorem}

\begin{advdefinition}{温度、压强与化学势}\label{advA2:def:intensives}
对单组分简单体系定义
\[
 \left(\frac{\partial S}{\partial U}\right)_{V,N}=\frac1T,
 \quad
 \left(\frac{\partial S}{\partial V}\right)_{U,N}=\frac pT,
 \quad
 \left(\frac{\partial S}{\partial N}\right)_{U,V}=-\frac\mu T.
\]
因此
\begin{equation}\label{advA2:eq:entropy-fundamental}
 \dd S=\frac{\dd U}{T}+\frac pT\dd V-\frac\mu T\dd N,
 \qquad
 \dd U=T\dd S-p\dd V+\mu\dd N.
\end{equation}
\end{advdefinition}
式\eqref{advA2:eq:entropy-fundamental}立即解释三个自发交换方向：热量从高温流向低温，机械平衡要求两侧压强相等，可交换粒子时平衡要求化学势相等。

\begin{advexamplebox}{小压强一侧为何仍可能膨胀}\label{advA2:ex:pressure-temperature}
两个子体系满足 $p_1<p_2$ 且 $T_1<T_2$。在固定各子体系内能与粒子数的熵虚变分中，解释为什么体积从体系2转给体系1仍可能使总熵增加；说明这是否就是自由绝热活塞的实际运动判据。
\end{advexamplebox}
\begin{advsolution}
令 $\dd V_1=-\dd V_2>0$，并按题设在状态空间中固定 $U_1,U_2,N_1,N_2$，则
\[
 \dd S_{\rm tot}
 =\left(\frac{p_1}{T_1}-\frac{p_2}{T_2}\right)\dd V_1.
\]
决定方向的是 $p/T$，不是单独的 $p$。例如取
$p_1=1,T_1=1,p_2=2,T_2=3$，虽然 $p_1<p_2$，却有
$p_1/T_1>p_2/T_2$，所以 $\dd V_1>0$ 对应熵增。单位检查：$p\dd V/T$ 的单位为能量除以温度，正是熵。这里比较的是受约束的邻近平衡态。实际绝热活塞移动时两侧一般还发生压强功引起的内能变化；总容器绝热并不意味着各侧 $\dd U_i=0$。因此不能据此断言自由活塞会自发向高压侧移动。
\end{advsolution}

\newcommand{\ASecondThermodynamicsLater}{%
\subsection{热力学函数}

\begin{advdefinition}{五个常用热力学势}\label{advA2:def:potentials}
对单组分体系，常用势及其自然变量为
\[
\begin{array}{c|c|c}
\text{势}&\text{定义}&\text{全微分}\\ \hline
U(S,V,N)&U&\dd U=T\dd S-p\dd V+\mu\dd N\\
H(S,p,N)&U+pV&\dd H=T\dd S+V\dd p+\mu\dd N\\
F(T,V,N)&U-TS&\dd F=-S\dd T-p\dd V+\mu\dd N\\
G(T,p,N)&U-TS+pV&\dd G=-S\dd T+V\dd p+\mu\dd N\\
\Omega(T,V,\mu)&F-\mu N&\dd\Omega=-S\dd T-p\dd V-N\dd\mu
\end{array}
\]
均匀宏观体系有 $\Omega=-pV$。源资料把 $-\Omega=pV$ 记作巨势；本卷采用统计物理中更常见的符号约定，并在后文保持一致。
\end{advdefinition}

\begin{advtheorem}{麦克斯韦关系}\label{advA2:thm:maxwell}
由各热力学势的混合偏导可交换，得到
\[
 \left(\frac{\partial T}{\partial V}\right)_S
 =-\left(\frac{\partial p}{\partial S}\right)_V,
 \quad
 \left(\frac{\partial T}{\partial p}\right)_S
 =\left(\frac{\partial V}{\partial S}\right)_p,
\]
\[
 \left(\frac{\partial S}{\partial V}\right)_T
 =\left(\frac{\partial p}{\partial T}\right)_V,
 \quad
 \left(\frac{\partial S}{\partial p}\right)_T
 =-\left(\frac{\partial V}{\partial T}\right)_p.
\]
\end{advtheorem}
例如最后两式分别来自 $F_{TV}=F_{VT}$ 与 $G_{Tp}=G_{pT}$。这一方法同时说明每个偏导中必须固定哪个自然变量，不能只凭符号记忆。

\begin{advexamplebox}{吉布斯--亥姆霍兹方程}\label{advA2:ex:gibbs-helmholtz}
证明
\[
 \boxed{\left[\frac{\partial(G/T)}{\partial T}\right]_p=-\frac{H}{T^2}}.
\]
\end{advexamplebox}
\begin{advsolution}
在定压下 $G_T=-S$，而 $G=H-TS$。因此
\[
 \left[\frac{\partial(G/T)}{\partial T}\right]_p
 =\frac{TG_T-G}{T^2}
 =\frac{-TS-(H-TS)}{T^2}=-\frac H{T^2}.
\]
两边单位均为熵除以温度。
\end{advsolution}

\begin{advexamplebox}{液膜表面张力与表面热力学}\label{advA2:ex:surface-thermodynamics}
面积为 $A$ 的液膜，其表面自由能为 $F_s=\gamma(T)A$。
\begin{enumerate}
 \item 求表面熵与表面内能。
 \item 若 $\gamma=C(T_c-T)$，等温增加面积 $\Delta A$，求外力功与吸热。
 \item 若 $\gamma=CT_c(1-T/T_c)^a$ 且 $1<a<2$，保持面积不变加热时表面层吸热还是放热？这是否意味着整个液体具有负热容？
\end{enumerate}
\end{advexamplebox}
\begin{advsolution}
表面功项为 $\gamma\dd A$，故
\[
 \dd F_s=-S_s\dd T+\gamma\dd A,
 \quad
 \boxed{S_s=-A\gamma'(T)},
 \quad
 \boxed{U_s=F_s+TS_s=A[\gamma-T\gamma']}.
\]
等温拉伸时
\[
 \boxed{W_{\rm ext}=\gamma\Delta A},
 \qquad
 \boxed{Q=T\Delta S_s=-T\gamma'\Delta A}.
\]
对线性定律，$W_{\rm ext}=C(T_c-T)\Delta A$，$Q=CT\Delta A>0$。

定面积时表面热容
\[
 C_A=T\left(\frac{\partial S_s}{\partial T}\right)_A
 =-TA\gamma''(T).
\]
当 $1<a<2$ 时 $\gamma''>0$，故表面\emph{贡献}为负，即升温时表面层向体相释放能量。这不等于完整孤立液体的总热容为负：表面层不是可独立隔离的宏观体系，体相热容与表面--体相能量交换必须同时计入。临界点附近还要考虑涨落，经验幂律也不能任意外推。
\end{advsolution}

\begin{advexamplebox}{重力场中理想气体与悬挂球争议}\label{advA2:ex:gravity-heat-capacity}
质量为 $m$ 的经典分子在 $0\le z\le H$ 内运动，外场势能为 $mgz$。除竖直势能外，每个分子的二次型自由度给出热容 $fk/2$。
\begin{enumerate}
 \item 求 $N$ 个分子的 $U(g,T)$ 与定 $g$ 热容 $C_g$，并比较 $g=0$ 与 $g\to\infty$。
 \item 对一般处在外场 $g$ 中的物体，令总质量为 $M$、质心高度为 $z(T,g)$，定义
 $\alpha_z=z^{-1}(\partial z/\partial T)_g$。推导 $(\partial C_g/\partial g)_T$。
 \item 地面上承托的球与天花板下悬挂的球吸收相同热量时，能否仅凭“膨胀时重力做功的正负”断言谁升温更多？
\end{enumerate}
\end{advexamplebox}
\begin{advsolution}
单粒子的竖直构型配分因子为
\[
 Z_z=\int_0^H e^{-mgz/(kT)}\dd z
 =\frac{1-e^{-x}}{\beta mg},
 \qquad x=\frac{mgH}{kT}.
\]
因此平均重力势能为
\[
 \langle mgz\rangle=-\frac{\partial\ln Z_z}{\partial\beta}
 =kT\left(1-\frac{x}{e^x-1}\right).
\]
于是
\[
 \boxed{U=N\left[\frac f2kT+kT\left(1-\frac{x}{e^x-1}\right)\right]},
\]
\[
 \boxed{C_g=Nk\left[\frac f2+1-\frac{x^2e^x}{(e^x-1)^2}\right]}.
\]
$x\to0$ 时括号中的重力贡献趋于零；$x\to\infty$ 时趋于1，所以
$C_g\to Nfk/2$ 与 $Nk(f/2+1)$。极限连续且热容为正。

把重力势能包括在体系能量中，亥姆霍兹自由能满足
\[
 \dd F=-S\dd T+Mz\dd g.
\]
混合偏导给出
\[
 \left(\frac{\partial S}{\partial g}\right)_T
 =-M\left(\frac{\partial z}{\partial T}\right)_g=-Mz\alpha_z.
\]
再对 $C_g=T(\partial S/\partial T)_g$ 求导：
\[
 \boxed{\left(\frac{\partial C_g}{\partial g}\right)_T
 =-MT\frac{\partial(z\alpha_z)}{\partial T}
 =-MT\left(\frac{\partial^2z}{\partial T^2}\right)_g}.
\]
直接把上面的气体质心 $z=(kT/mg)-H/(e^x-1)$ 代入，可复现 $C_g$ 对 $g$ 的变化。

对两个球，支承方式改变了应力分布、质心高度函数和热容。升温量应由完整的
$\Delta T=Q/C_g$ 决定，不能把“内能增加”与“温度增加”直接等同。若假设仅靠支承方向便能无条件得到固定温差，把两球交替接触热库并互换支承即可构造从单一热源循环取功的装置，违背第二定律。因此争议推理中不严谨的是忽略状态方程、弹性能和热容，仅按重力功符号比较温升。

\begin{center}\begin{minipage}{\linewidth}\centering
\begin{tikzpicture}[scale=.9]
 \begin{scope}[xshift=-2.7cm]
  \draw[ground] (-1.2,0)--(1.2,0);
  \draw[very thick,fill=DeepBlue!6] (0,.72) circle(.72);
  \draw[vector] (0,.72)--(0,.15) node[midway,right]{$Mg$};
  \node at (0,-.45) {地面承托};
 \end{scope}
 \begin{scope}[xshift=2.7cm]
  \draw[ground] (-1.2,2.25)--(1.2,2.25);
  \draw[very thick] (0,2.25)--(0,1.45);
  \draw[very thick,fill=DeepBlue!6] (0,.73) circle(.72);
  \draw[vector] (0,.73)--(0,.16) node[midway,right]{$Mg$};
  \node at (0,-.45) {绳悬挂};
 \end{scope}
\end{tikzpicture}
\bookcaption{支承方式改变约束和应力，不能只按重力功的正负判断温升}\label{advA2:fig:supported-balls}
\end{minipage}\end{center}
图\ref{advA2:fig:supported-balls}只表达约束差别；比较温升时必须为两种约束分别求 $z(T,g)$ 与 $C_g$。
\end{advsolution}

\subsection{相平衡与相变}

同一温度、压强下，相共存要求两相的化学势相等。对单组分物质，$\dd\mu=-s\dd T+v\dd p$，故沿共存线
\[
 \frac{\dd p}{\dd T}=\frac{s_2-s_1}{v_2-v_1}
 =\frac{\ell}{T(v_2-v_1)}.
\]
这里 $s,v,\ell$ 必须统一为每摩尔量或每单位质量量。分母很小时不能直接使用“液体体积可忽略”的蒸气近似。

\paragraph{由自由能得到等面积条件。}
等温线上若状态方程出现回折，机械稳定要求 $(\partial p/\partial v)_T<0$；斜率为正的部分不是可稳定保持的均匀相。
记摩尔 Helmholtz 自由能为 $f$，则 $\dd f=-p\dd v$。两端相同 $p_*$ 且 $\mu=f+p_*v$ 相等，给出
\[
 \int_{v_1}^{v_2}[p(v,T)-p_*]\dd v=0.
\]
这就是 Maxwell 等面积条件，确定的是两相共存端点；$(\partial p/\partial v)_T=0$ 给出的旋节边界只标记均匀相局部稳定性丧失。临界点处两端合并，常见光滑状态方程满足一、二阶体积导数同时为零。解题时应先问所求是共存、局部稳定，还是临界条件，三者不能互换。


\begin{advtheorem}{范德瓦耳斯流体的临界点与对比状态方程}\label{advA2:thm:vdw}
对
\[
 \left(p+\frac{N^2a}{V^2}\right)(V-Nb)=NkT,
\]
临界常数和对比状态方程为
\[
 \boxed{V_c=3Nb,\quad p_c=\frac{a}{27b^2},\quad kT_c=\frac{8a}{27b}},
\]
\[
 \boxed{\left(\pi+\frac3{v^2}\right)(3v-1)=8t},
 \quad v=\frac V{V_c},\ \pi=\frac p{p_c},\ t=\frac T{T_c}.
\]
\end{advtheorem}
证明：把压强写成
$p=NkT/(V-Nb)-N^2a/V^2$。临界等温线的水平拐点满足
$(\partial p/\partial V)_T=(\partial^2p/\partial V^2)_T=0$；联立即得三项临界常数。代入无量纲变量后，材料参数 $a,b$ 全部消失，得到对比状态方程。

两相平衡要求 $T_1=T_2,p_1=p_2,\mu_1=\mu_2$。一级相变中 $\mu$ 的一阶导数，即摩尔熵和摩尔体积，可以跃变；连续相变的一阶导数连续而二阶导数可跃变。

\begin{advtheorem}{克拉珀龙与埃伦费斯关系}\label{advA2:thm:clapeyron}
一级相变共存曲线满足
\[
 \boxed{\frac{\dd p}{\dd T}=\frac{\Delta s}{\Delta v}=\frac{L}{T\Delta v}}.
\]
连续相变在相界上若 $s,v$ 连续，则
\[
 \boxed{\frac{\dd p}{\dd T}=\frac{\Delta c_p}{Tv\Delta\alpha}
 =\frac{\Delta\alpha}{\Delta\kappa_T}}.
\]
\end{advtheorem}
证明：两相化学势相等，沿共存线微分
$-s_1\dd T+v_1\dd p=-s_2\dd T+v_2\dd p$，即得克拉珀龙式。连续相变时再分别对 $\Delta s=0$、$\Delta v=0$ 沿相界求全微分，利用
$s_T=c_p/T,s_p=-v\alpha,v_T=v\alpha,v_p=-v\kappa_T$，便得两条埃伦费斯关系。

\begin{advexamplebox}{超导临界场、潜热与热容突变}\label{advA2:ex:superconducting-transition}
正常--超导相变的热力学临界场满足
\[
 H_c(T)=H_0\left(1-\frac{T^2}{T_c^2}\right).
\]
求相变潜热，并求零场二级相变点的单位体积热容突变。
\end{advexamplebox}
\begin{advsolution}
采用SI制，凝聚能密度满足
\[
 g_n-g_s=\frac{\mu_0H_c^2}{2}.
\]
因此
\[
 s_n-s_s=-\mu_0H_c\frac{\dd H_c}{\dd T},
 \qquad
 L=T(s_n-s_s)=-\mu_0TH_cH_c'.
\]
代入给定函数，
\[
 \boxed{L(T)=2\mu_0H_0^2\frac{T^2}{T_c^2}
 \left(1-\frac{T^2}{T_c^2}\right)}.
\]
在零场端点 $T=T_c$，$H_c=0$，所以潜热为零；但
\[
 c_s-c_n=T\frac{\dd(s_s-s_n)}{\dd T}
 =\mu_0T[(H_c')^2+H_cH_c'']
\]
给出
\[
 \boxed{\Delta c\big|_{T_c}=\frac{4\mu_0H_0^2}{T_c}}.
\]
若采用高斯制，只需把 $\mu_0H_c^2/2$ 换成 $H_c^2/(8\pi)$；物理结论不变。
\end{advsolution}

平均场海森堡模型把每个自旋受到的近邻交换作用替换成与平均磁化强度成正比的分子场。对自旋 $1/2$，令单自旋磁矩 $\mu_s=\hbar\gamma/2$，饱和磁矩 $M_0=N\mu_s$。若每个自旋有 $n$ 个近邻，交换能标为 $\hbar J$，零外场自洽方程可写成
\begin{equation}\label{advA2:eq:weiss-self-consistent}
 m=\tanh\!\left(\frac{T_c}{T}m\right),
 \qquad m=\frac M{M_0},\qquad
 kT_c=\frac{n\hbar J}{4}.
\end{equation}

\begin{advexamplebox}{外斯模型的临界温度}\label{advA2:ex:weiss-critical}
由自洽方程在正温度下，证明自发磁化的非零解只在 $T<T_c=n\hbar J/(4k)$ 时出现，并检查 $T=T_c$。
\end{advexamplebox}
\begin{advsolution}
$m=0$ 总是解。临界点附近展开 $\tanh x=x-x^3/3+\cdots$：
\[
 m=\frac{T_c}{T}m-\frac13\left(\frac{T_c}{T}m\right)^3+\cdots.
\]
非零解要求线性项斜率 $T_c/T>1$，即 $T<T_c$。在 $T=T_c$，由于 $|\tanh m|<|m|$ 对所有 $m\ne0$ 成立，仍只有 $m=0$。$T>T_c$ 时右端曲线在原点斜率小于1，且始终被 $m$ 直线压住，故没有别的交点。
\end{advsolution}

\begin{advexamplebox}{铁磁相变曲线与临界指数}\label{advA2:ex:weiss-curve}
证明零场相变曲线可写成
\[
 \frac{T}{T_c}=\frac{2m}{\ln[(1+m)/(1-m)]},
\]
求 $M_0$ 并确定 $T\to T_c^-$ 时磁化强度的临界指数。
\end{advexamplebox}
\begin{advsolution}
对式\eqref{advA2:eq:weiss-self-consistent}取反双曲正切：
\[
 \frac{T_c}{T}m=\operatorname{artanh}m
 =\frac12\ln\frac{1+m}{1-m},
\]
即得题式。自旋全同向时
$\boxed{M_0=N\hbar\gamma/2}$。临界附近由上例展开式，
\[
 m^2\simeq3\left(1-\frac T{T_c}\right),
 \quad
 \boxed{M\propto(T_c-T)^{1/2}}.
\]
平均场磁化临界指数为 $1/2$。
\begin{center}\begin{minipage}{\linewidth}\centering
\begin{tikzpicture}[scale=.9]
 \draw[axis] (0,0)--(4.5,0) node[right]{$T/T_c$};
 \draw[axis] (0,0)--(0,2.8) node[above]{$M/M_0$};
 \draw[curve] plot[smooth] coordinates {
   (0,2.4) (.72,2.40) (1.44,2.37) (2.16,2.18)
   (2.88,1.70) (3.24,1.26) (3.42,.91) (3.60,0)};
 \draw[guide] (3.6,0)--(3.6,2.55);
 \node[below] at (3.6,0) {$1$};
 \node[left] at (0,2.4) {$1$};
\end{tikzpicture}
\bookcaption{平均场磁化强度在临界点连续降为零，近临界区呈平方根律}\label{advA2:fig:weiss-curve}
\end{minipage}\end{center}
\end{advsolution}

\begin{advexamplebox}{居里--外斯定律}\label{advA2:ex:curie-weiss}
在线性弱场区求磁化率，并证明 $1/\chi$ 与 $T-T_c$ 成正比。
\end{advexamplebox}
\begin{advsolution}
把分子场写成 $B_{\rm eff}=B+\lambda M$。弱场高温下
\[
 M=N\mu_s\tanh\frac{\mu_s(B+\lambda M)}{kT}
 \simeq\frac{N\mu_s^2}{kT}(B+\lambda M).
\]
定义 $T_c=N\mu_s^2\lambda/k$，解得
\[
 \boxed{\chi=\frac{\partial M}{\partial B}
 =\frac{N\mu_s^2/k}{T-T_c}},
 \qquad
 \boxed{\frac1\chi=\frac{k}{N\mu_s^2}(T-T_c)}.
\]
当 $\lambda\to0$ 时退化为居里定律 $\chi\propto1/T$。
\end{advsolution}

朗道理论用序参量 $\phi$ 的对称性组织相变。若体系在 $\phi\to-\phi$ 下不变，最低阶展开可写成
\[
 F=F_0+a(T-T_c)\phi^2+b\phi^4+\cdots,
 \qquad b>0.
\]
$T<T_c$ 时平衡序参量为
$\phi_\pm=\pm\sqrt{-a(T-T_c)/(2b)}$。熵与热容分别为
\[
 S=-\left(\frac{\partial F}{\partial T}\right)_\phi,
 \qquad C=T\left(\frac{\partial S}{\partial T}\right),
\]
从而在平均场近似中热容出现有限跃变
$\Delta C=a^2T_c/(2b)$。

\begin{advexamplebox}{六次朗道势与一级相变}\label{advA2:ex:landau-first-order}
证明当
\[
 F(\phi)=F_0+a\phi^2-b\phi^4+c\phi^6,
 \qquad a,b,c>0,
\]
随控制参量改变时可以发生一级相变，并求两相等自由能时的非零序参量与 $a$。
\end{advexamplebox}
\begin{advsolution}
令 $x=\phi^2\ge0$。非零驻点满足
\[
 a-2bx+3cx^2=0.
\]
一级相变共存时，$\phi=0$ 与非零极小值具有相同自由能，故还需
$ax-bx^2+cx^3=0$。两式联立得
\[
 \boxed{x=\phi^2=\frac b{2c}},
 \qquad
 \boxed{a=\frac{b^2}{4c}}.
\]
在共存点序参量从0跳到 $\sqrt{b/(2c)}$，故是一级相变；$c>0$ 保证大 $|\phi|$ 时自由能有下界。若加入空间非均匀与电磁耦合，得到朗道--金兹堡泛函
\[
 F=F_n+\alpha|\psi|^2+\frac\beta2|\psi|^4
 +\frac1{2m}\left|(-i\hbar\nabla-2e\vect A)\psi\right|^2
 +\frac{|\vect B|^2}{2\mu_0}.
\]
\end{advsolution}

}

\begin{advexamplebox}{欧拉齐次函数定理}\label{advA2:ex:euler-homogeneous}
若 $A(\lambda x_1,\ldots,\lambda x_n)=\lambda^kA(x_1,\ldots,x_n)$，证明
\[
 \boxed{kA=\sum_i x_i\frac{\partial A}{\partial x_i}}.
\]
\end{advexamplebox}
\begin{advsolution}
固定 $x_i$，对缩放参数求导：
\[
 \frac{\dd}{\dd\lambda}A(\lambda x_i)
 =\sum_i x_i\frac{\partial A}{\partial(\lambda x_i)}
 =k\lambda^{k-1}A(x_i).
\]
令 $\lambda=1$ 即得结论。对广延热力学量 $k=1$；对强度量 $k=0$，相应加权和为零。
\end{advsolution}

\begin{advexamplebox}{萨克--泰特洛德方程给出的理想气体热力学}\label{advA2:ex:sackur-tetrode}
单原子经典理想气体满足
\[
 \frac{S}{Nk}=\ln\!\left[
 \frac VN\left(\frac{4\pi m}{3h^2}\frac UN\right)^{3/2}
 \right]+\frac52.
\]
以 $N,V,T$ 为自变量求 $U,S,p,\mu$。
\end{advexamplebox}
\begin{advsolution}
由熵基本关系，
\[
 \frac1T=\left(\frac{\partial S}{\partial U}\right)_{V,N}
 =\frac{3Nk}{2U},
\]
故
\[
 \boxed{U=\frac32NkT},\qquad
 \boxed{p=T\left(\frac{\partial S}{\partial V}\right)_{U,N}
 =\frac{NkT}{V}}.
\]
代回熵式，
\[
 \boxed{S=Nk\left\{\ln\!\left[
 \frac VN\left(\frac{2\pi mkT}{h^2}\right)^{3/2}\right]+\frac52\right\}}.
\]
再由 $\mu=-T(\partial S/\partial N)_{U,V}$，或用下述欧拉关系，得到
\[
 \boxed{\mu=-kT\ln\!\left[
 \frac VN\left(\frac{2\pi mkT}{h^2}\right)^{3/2}\right]}.
\]
$U$ 与 $pV$ 均与 $NkT$ 同量纲；稀薄极限 $V/N\to\infty$ 时 $\mu\to-\infty$，符合经典气体极限。
\end{advsolution}

熵式和能量式分别称为熵表象与能量表象的基本关系。绝对零点附近还需第三定律约束积分常数。

\begin{advtheorem}{热力学第三定律的能量表述}\label{advA2:thm:third-law}
在基态不具有宏观简并的通常情形，$T\to0$ 时熵趋于与外参量无关的常数，可取为零，且
\[
 \lim_{S\to0}\left(\frac{\partial U}{\partial S}\right)_{V,N}=0,
 \qquad
 U(S,V,N)=U_0(V,N)+\int_0^S T\,\dd S.
\]
\end{advtheorem}
广延性还给出 $U(S,V,N)=Nu(S/N,V/N)$，以后用小写字母表示每粒子量。

\begin{advexamplebox}{化学势关系与吉布斯--杜亥姆关系}\label{advA2:ex:gibbs-duhem}
证明
\[
 \mu=u-Ts+pv,
 \qquad U=TS-pV+\mu N,
\]
并写出吉布斯--杜亥姆关系。
\end{advexamplebox}
\begin{advsolution}
$U(S,V,N)$ 是一次齐次函数。对它应用例\ref{advA2:ex:euler-homogeneous}，再用
$U_S=T,U_V=-p,U_N=\mu$，得
\[
 \boxed{U=TS-pV+\mu N}.
\]
除以 $N$ 即有 $\mu=u-Ts+pv$。对欧拉式求微分，并与
$\dd U=T\dd S-p\dd V+\mu\dd N$ 比较，消去同类项：
\[
 \boxed{S\dd T-V\dd p+N\dd\mu=0}.
\]
这说明三个强度量不能任意独立改变。
\end{advsolution}

对定粒子数体系，焓 $H=U+pV$ 满足 $\dd H=T\dd S+V\dd p$。定义
\[
 C_p=T\left(\frac{\partial S}{\partial T}\right)_p
 =\left(\frac{\partial H}{\partial T}\right)_p,
 \quad
 C_V=T\left(\frac{\partial S}{\partial T}\right)_V
 =\left(\frac{\partial U}{\partial T}\right)_V,
\]
\[
 \alpha=\frac1V\left(\frac{\partial V}{\partial T}\right)_p,
 \quad
 \beta=\frac1p\left(\frac{\partial p}{\partial T}\right)_V,
 \quad
 \kappa_T=-\frac1V\left(\frac{\partial V}{\partial p}\right)_T,
 \quad
 \kappa_S=-\frac1V\left(\frac{\partial V}{\partial p}\right)_S.
\]

\begin{advtheorem}{能态定理}\label{advA2:thm:energy-state}
\[
 \boxed{\left(\frac{\partial U}{\partial V}\right)_T+p
 =T\left(\frac{\partial p}{\partial T}\right)_V}.
\]
\end{advtheorem}
证明：令 $F=U-TS$，则
$\dd F=-S\dd T-p\dd V$，因而
$(\partial S/\partial V)_T=(\partial p/\partial T)_V$。又
\[
 \left(\frac{\partial U}{\partial V}\right)_T
 =\left(\frac{\partial F}{\partial V}\right)_T
 +T\left(\frac{\partial S}{\partial V}\right)_T
 =-p+T\left(\frac{\partial p}{\partial T}\right)_V.
\]

\begin{advexamplebox}{膨胀系数、压强系数与等温压缩系数}\label{advA2:ex:alpha-beta-kappa}
证明 $\alpha=\beta\kappa_Tp$。
\end{advexamplebox}
\begin{advsolution}
把状态方程写成 $V=V(T,p)$。在定体积线上
\[
 0=\left(\frac{\partial V}{\partial T}\right)_p\dd T
 +\left(\frac{\partial V}{\partial p}\right)_T\dd p,
\]
所以
\[
 \left(\frac{\partial p}{\partial T}\right)_V
 =-\frac{V\alpha}{-V\kappa_T}=\frac{\alpha}{\kappa_T}.
\]
左侧又等于 $p\beta$，故 $\boxed{\alpha=p\beta\kappa_T}$。理想气体给出
$\alpha=1/T,\beta=1/T,\kappa_T=1/p$，校验成立。
\end{advsolution}

\begin{advexamplebox}{热容比与压缩系数比}\label{advA2:ex:cp-cv-kappa}
证明
\[
 \boxed{\frac{C_p}{C_V}=\frac{\kappa_T}{\kappa_S}}.
\]
\end{advexamplebox}
\begin{advsolution}
把熵写成 $S(T,V)$。由麦克斯韦关系
$(\partial S/\partial V)_T=(\partial p/\partial T)_V$，
\[
 \dd S=\frac{C_V}{T}\dd T
 +\left(\frac{\partial p}{\partial T}\right)_V\dd V.
\]
沿绝热线令 $\dd S=0$，再与状态方程微分联立，可得
\[
 \kappa_S=\kappa_T-\frac{TV\alpha^2}{C_p}.
\]
另一方面
$C_p-C_V=TV\alpha^2/\kappa_T$。消元便得
$\kappa_S=\kappa_TC_V/C_p$，即所求。稳定体系中
$C_p\ge C_V>0$，所以 $\kappa_T\ge\kappa_S>0$，方向也正确。
\end{advsolution}

\begin{advexamplebox}{两种热容之差}\label{advA2:ex:cp-minus-cv}
证明
\[
 \boxed{C_p-C_V=\alpha\beta pVT}.
\]
\end{advexamplebox}
\begin{advsolution}
由
\[
 C_p-C_V=T\left[\left(\frac{\partial S}{\partial T}\right)_p
 -\left(\frac{\partial S}{\partial T}\right)_V\right]
\]
及链式法则，
\[
 C_p-C_V
 =T\left(\frac{\partial S}{\partial V}\right)_T
 \left(\frac{\partial V}{\partial T}\right)_p
 =T(p\beta)(V\alpha).
\]
因此得到题式。用例\ref{advA2:ex:alpha-beta-kappa}也可写成
$C_p-C_V=TV\alpha^2/\kappa_T\ge0$；理想气体代入后为 $Nk$，与
$C_p=5Nk/2,C_V=3Nk/2$ 一致。
\end{advsolution}

\ASecondThermodynamicsLater
% AdvancedThermalQuality20261023-v1
% Editable statistical mechanics figures and the quantitative Bose heat curve.
\providecommand{\ThermalFigure}[3]{\begin{qualityfigure}\begin{tikzpicture}[>=Latex,font=\small,x=.8cm,y=.8cm]#1\end{tikzpicture}\figcap{#2}\label{#3}\end{qualityfigure}}
\newcommand{\AQMicrocanonical}{\ThermalFigure{
 \draw[->](-2.5,0)--(2.5,0)node[right]{$q$};\draw[->](0,-1.7)--(0,1.7)node[above]{$p$};
 \draw[DeepBlue,thick](0,0)ellipse(2 and 1.2);\draw[DeepBlue,thick](0,0)ellipse(1.7 and 1);\node[right]at(3,1){能壳上的相空间体积 $\omega$};
}{只在给定能壳上归一化概率密度；壳外不属于本次熵极值问题。}{advA2:quality:micro-q}}
\newcommand{\AAMicrocanonical}{\ThermalFigure{
 \foreach\x in {0,1,2,3}{\draw[fill=DeepBlue!12](\x,0)rectangle(\x+.8,1.2);}
 \foreach\x/\h in {6/.4,7/2,8/1.6,9/.8}{\draw[fill=TextGray!15](\x,0)rectangle(\x+.8,\h);}
 \node[below]at(1.9,-.4){等体积小胞：均匀密度};\node[below]at(7.9,-.4){同一归一化：非均匀密度};
}{把同一能壳分成等体积小胞，两组柱面积相同，仅分布不同。}{advA2:quality:micro-a}
图~\ref{advA2:quality:micro-a}比较的是同一能壳和同一测度下的密度；若改变能壳体积，就不是同一个极值问题。相对熵非负可由 $x\ln x-x+1\ge0$ 对均匀参考测度积分得到，等号要求密度比几乎处处等于1，因而也给出最大值的唯一性。
}
\newcommand{\AQGrand}{\ThermalFigure{
 \draw[thick](0,0)rectangle(4,2);\node at(2,1){系统，$V$ 固定};\node[draw](R)at(8,1){巨大热、粒子库 $T,\mu$};
 \draw[<->,DeepBlue,thick](4,1.5)--(6,1.5)node[midway,above]{能量};\draw[<->,thick](4,.5)--(6,.5)node[midway,below]{粒子};
}{巨正则体系允许能量和粒子数涨落，但体积作为外控参数固定。}{advA2:quality:grand-q}}
\newcommand{\AAGrand}{\ThermalFigure{
 \node[draw](X)at(4,1){$\ln\Xi(\mu,V,T)$};\node[draw](N)at(0,0){$\partial_\mu\ln\Xi$};\node[draw](P)at(4,-1){$\partial_V\ln\Xi$};\node[draw](S)at(8,0){$\partial_T\ln\Xi$};
 \draw[->](X)--(N);\draw[->](X)--(P);\draw[->](X)--(S);\node[above]at(4,2){每次偏导都须固定其余自然变量};
}{粒子数、压强、熵来自同一巨配分函数的不同参数响应。}{advA2:quality:grand-a}
图~\ref{advA2:quality:grand-a}中压强的一般公式是 $p=kT(\partial\ln\Xi/\partial V)_{T,\mu}$。只有均匀、广延的体相体系，忽略边界修正并满足 $\Omega=-pV$ 时，才可进一步写成 $p=(kT/V)\ln\Xi$。固定逸度求温度导数与固定化学势求导不同，因为 $z=e^{\beta\mu}$ 本身含温度。
}
\newcommand{\AQDensityStates}{\ThermalFigure{
 \draw[->](-2,0)--(2,0)node[right]{$p$};\draw[very thick](-1.2,0)--(1.2,0);\node[below]at(0,-.5){一维动量区间};
 \begin{scope}[shift={(5,0)}]\draw[thick](0,0)circle(1.3);\draw[DeepBlue](0,0)circle(1.1);\node[below]at(0,-1.7){二维动量圆壳};\end{scope}
 \begin{scope}[shift={(10,0)}]\draw[thick](0,0)circle(1.3);\draw[dashed](0,0)ellipse(1.3 and .4);\draw[->](0,0)--(45:1.3)node[right]{$p$};\node[below]at(0,-1.7){三维动量球壳};\end{scope}
}{先数能量不超过给定值的动量态，再由壳体积求能量态密度。}{advA2:quality:dos-q}}
\newcommand{\AADensityStates}{\ThermalFigure{
 \draw[->](0,0)--(4,0)node[right]{$\varepsilon$};\draw[->](0,0)--(0,3)node[above]{$g$};
 \draw[DeepBlue,domain=.25:3.5,samples=60]plot(\x,{1/sqrt(\x)});\draw[thick](.25,1)--(3.5,1);\draw[TextGray,domain=.25:3.5,samples=60]plot(\x,{sqrt(\x)});
 \node[right,DeepBlue]at(3.5,.53){$d=1$};\node[right]at(3.5,1){$d=2$};\node[right,TextGray]at(3.5,1.87){$d=3$};
 \node[above]at(2,3){非相对论：指数 $d/2-1$};
 \begin{scope}[shift={(7,0)}]\draw[->](0,0)--(4,0)node[right]{$\varepsilon$};\draw[->](0,0)--(0,3)node[above]{$g$};\draw[thick](.1,.5)--(3,.5);\draw[DeepBlue](0,0)--(3,1.5);\draw[TextGray,domain=0:3,samples=60]plot(\x,{.25*\x*\x});\node[above]at(2,3){无质量：指数 $d-1$};\node[right]at(3,.5){$d=1$};\node[right,DeepBlue]at(3,1.5){$d=2$};\node[right,TextGray]at(3,2.25){$d=3$};\end{scope}
}{两种色散均按 $d=1,2,3$ 比较幂律，曲线归一化系数仅为示意。}{advA2:quality:dos-a}
图~\ref{advA2:quality:dos-a}中的幂指数由两个因子合成：动量球壳给 $p^{d-1}\dd p$，从动量改成能量还要乘 $\dd p/\dd\varepsilon$。所以非相对论得到 $\varepsilon^{d/2-1}$，无质量粒子得到 $\varepsilon^{d-1}$；直接把球壳半径改名为能量会漏掉雅可比。
}
\newcommand{\AQFreePartition}{\ThermalFigure{
 \draw[->](0,0)--(4,0)node[right]{$p$};\draw[->](0,0)--(0,3)node[above]{$\varepsilon$};\draw[DeepBlue,thick,domain=0:3,samples=50]plot(\x,{.25*\x*\x});\node[right]at(4,1.5){$p^2/(2m)$};
 \begin{scope}[shift={(7,0)}]\draw[->](0,0)--(4,0)node[right]{$p$};\draw[->](0,0)--(0,3)node[above]{$\varepsilon$};\draw[DeepBlue,thick](0,0)--(3,2.4);\node[right]at(3,2.4){$pc$};\end{scope}
}{两种单粒子色散关系；三维空间和统计权重保持相同。}{advA2:quality:free-q}}
\newcommand{\AAFreePartition}{\ThermalFigure{
 \draw[->](0,0)--(6,0)node[right]{$\beta\varepsilon$};\draw[->](0,0)--(0,3)node[above]{积分权重};
 \draw[DeepBlue,thick,domain=.001:5.5,samples=90]plot(\x,{3*sqrt(\x)*exp(-\x)});
 \draw[thick,domain=0:5.5,samples=90]plot(\x,{2*\x*\x*exp(-\x)});
 \node[right,align=left]at(6.5,1.6){非相对论：$\varepsilon^{1/2}e^{-\beta\varepsilon}$\\无质量：$\varepsilon^2e^{-\beta\varepsilon}$};
}{配分积分把态密度与玻尔兹曼权重相乘，能量平均再多乘一因子 $\varepsilon$。}{advA2:quality:free-a}
图~\ref{advA2:quality:free-a}中两权重的幂次不同。令 $u=\beta\varepsilon$ 后，积分的温度依赖全部提出为 $T^{3/2}$ 或 $T^3$，而无量纲积分为常数；对 $\ln z_1$ 求温度或 $\beta$ 导数便得到对应均能。对这里的固定粒子数模型，两者均有 $pV=NkT$。
}
\newcommand{\AQClassicalEntropy}{\ThermalFigure{
 \draw[thick](0,0)rectangle(5,2);\foreach\x/\y in {1/.6,2.7/1.4,4/.5}{\fill[DeepBlue](\x,\y)circle(3pt);}
 \draw[<->,thick](1,.9)..controls(1.8,2.4)and(2.8,2.4)..(2.7,1.7);\node[right]at(6,1){交换全同粒子的标号};
}{同一物理构形可被不同粒子标号重复记录，经典全同粒子配分函数须处理这一重复。}{advA2:quality:classical-s-q}}
\newcommand{\AAClassicalEntropy}{\ThermalFigure{
 \node[draw,align=center](L)at(2,1){三粒子标号的六种排列\\$123,132,213,231,312,321$};\node[draw](R)at(9,1){一个无标号构形};\draw[->,thick](L)--node[above]{除以 $3!$}(R);
}{$N!$ 修正将标号重复还原为物理态计数，而不是删除真实可区分的构型。}{advA2:quality:classical-s-a}
图~\ref{advA2:quality:classical-s-a}解释 $1/N!$ 的来源。用斯特林公式 $\ln N!\simeq N\ln N-N$ 后，$S$ 才具有热力学极限下的广延性；有限 $N$ 仍有次广延修正。展开成 $\ln T$、$\ln(V/N)$ 的写法只是选定单位后的缩写，严谨的量纲检查应使用完整的无量纲组合 $V/(N\lambda_T^3)$。
}
\newcommand{\AQOscillators}{\ThermalFigure{
 \foreach\s/\i in {0/1,4/2,8/N}{\begin{scope}[shift={(\s,0)}]\draw[thick](0,0)--(0,1.2);\draw[decorate,decoration={coil,segment length=2mm,amplitude=1mm},thick](0,.6)--(1.7,.6);\draw[thick](1.7,.1)rectangle(2.5,1.1);\node at(2.1,.6){$\i$};\end{scope}}
 \node[above]at(5,2){固定在不同位置、彼此独立的振子，频率均为 $\omega$};
}{振子由不同位置标识，统计计数不能照搬自由全同粒子的 $1/N!$。}{advA2:quality:osc-q}}
\newcommand{\AAOscillators}{\ThermalFigure{
 \draw[->](-2.5,0)--(2.5,0)node[right]{$q$};\draw[->](0,-1.8)--(0,1.8)node[above]{$p$};\draw[DeepBlue,thick](0,0)ellipse(2 and 1.2);
 \node[right,align=left]at(3,1){$H=p^2/(2m)+m\omega^2q^2/2$\\等能线为椭圆\\两个独立二次项各贡献 $kT/2$};
}{相空间高斯积分在 $q,p$ 两方向因子化，再对独立振子取乘积。}{advA2:quality:osc-a}
图~\ref{advA2:quality:osc-a}中的椭圆不是位置空间轨道，而是单振子的相空间等能线。两维高斯积分分别产生 $T^{1/2}$，故 $z\propto T$；$N$ 个有固定标识的位置对应 $Z=z^N$，而不是再除以 $N!$。
}
\newcommand{\AQRotors}{\ThermalFigure{
 \draw[fill=DeepBlue!10,thick](0,0)ellipse(1.8 and 1);\foreach\x/\y/\lab in {3/0/1,0/2/2,-1.8/-1.5/3}{\draw[->,thick](0,0)--(\x,\y)node[right]{$I_{\lab}$ 主轴};}
 \node[right]at(4.5,1){非线性刚体，三个主惯量均非零};
}{主轴是固定于刚体的坐标方向，取向还需用整个旋转群计数。}{advA2:quality:rotor-q}}
\newcommand{\AARotors}{\ThermalFigure{
 \draw[thick](0,0)ellipse(2 and 1);\draw[dashed](0,0)ellipse(2 and .35);\draw[->](0,0)--(3,0)node[right]{$L_1$};\draw[->](0,0)--(0,2)node[above]{$L_2$};\draw[->](0,0)--(-1.5,-1.6)node[left]{$L_3$};
 \node[right,align=left]at(4,1){$H=\sum_iL_i^2/(2I_i)$\\角动量积分给 $T^{3/2}\sqrt{I_1I_2I_3}$\\取向体积为 $8\pi^2/\sigma$};
}{角动量空间等能面为椭球，取向对称数另外消除同一取向的重复。}{advA2:quality:rotor-a}
图~\ref{advA2:quality:rotor-a}中三个主轴二次项各独立贡献一个高斯积分。取向对称数 $\sigma$ 处理分子自身的旋转等价，$N!$ 则处理不同全同分子交换，两者不是同一种重复，不能相互替代。线性分子的轴向惯量为零，必须改用两个转动自由度的模型。
}
\newcommand{\AQDipoleStatistics}{\ThermalFigure{
 \draw[->,thick](0,0)--(0,3)node[above]{$\vect E$};\draw[->,DeepBlue,thick](0,0)--(2,1.7)node[right]{$\vect d$};\draw(0,1)arc(90:40.4:1)node[right]{$\theta$};\node[right]at(4,1.5){只计取向能 $-dE\cos\theta$};
}{偶极取向相对外场的极角为 $\theta$，方位角仍需积分。}{advA2:quality:dip-stat-q}}
\newcommand{\AADipoleStatistics}{\ThermalFigure{
 \draw[thick](0,0)circle(1.5);\draw[DeepBlue,thick](-1.2,.9)--(1.2,.9);\draw[->](0,0)--(1.2,.9);\draw[->](0,0)--(0,2)node[above]{$\vect E$};
 \node[right,align=left]at(3,1){取向带测度 $\dd\Omega=2\pi\sin\theta\dd\theta$\\令 $u=\cos\theta$，积分上下限为 $-1,1$};
}{取向球上的纬带面积与玻尔兹曼权重相乘；角度并不是均匀分布在 $[0,\pi]$。}{advA2:quality:dip-stat-a}
图~\ref{advA2:quality:dip-stat-a}中的 $\sin\theta$ 来自球面面积元，不能漏掉。弱场展开时球面对称使 $\langle\cos\theta\rangle_0=0$、$\langle\cos^2\theta\rangle_0=1/3$，直接给出线性响应中的 $1/3$，也是朗之万函数展开的独立核验。
}
\newcommand{\AQQuantumDistributions}{\ThermalFigure{
 \foreach\s/\lab in {0/费米态,6/玻色态}{\begin{scope}[shift={(\s,0)}]\draw[thick](-1,0)--(2,0);\node[above]at(.5,.7){能级 $\varepsilon_i$，\lab};\end{scope}}
 \node[below]at(.5,-.6){占据数只能为 $0,1$};\node[below]at(6.5,-.6){占据数可为 $0,1,2,\ldots$};
}{两种统计的区别先体现在同一个单粒子态允许多少粒子占据。}{advA2:quality:quantum-q}}
\newcommand{\AAQuantumDistributions}{\ThermalFigure{
 \draw[->](0,0)--(5,0)node[right]{$x=(\varepsilon-\mu)/(kT)$};\draw[->](0,0)--(0,4)node[above]{$\langle n\rangle$};
 \draw[DeepBlue,thick,domain=0:4.5,samples=70]plot(\x,{1/(exp(\x)+1)});
 \draw[thick,domain=.25:4.5,samples=100]plot(\x,{1/(exp(\x)-1)});
 \draw[dashed,domain=0:4.5,samples=70]plot(\x,{exp(-\x)});\node[right,align=left]at(6,2){上：玻色\\中：经典极限 $e^{-x}$\\下：费米（图中 $x>0$）};
}{同一能级的平均占据；两种量子分布在 $x\gg1$ 时趋于同一经典尾部。}{advA2:quality:quantum-a}
图~\ref{advA2:quality:quantum-a}中玻色支仅在 $x>0$ 有有限正占据。对有限系统的单基态几何级数，严格收敛要求 $\mu<\varepsilon_0$；热力学极限中的凝聚态需单独处理，才可说 $\mu$ 从下方趋于基态能。费米支则自动限制每个态的占据不超过1。
}
\newcommand{\AQFermiCapacity}{\ThermalFigure{
 \draw[->](0,0)--(6,0)node[right]{$\varepsilon$};\draw[->](0,0)--(0,2.5)node[above]{$f$};\draw[thick](0,2)--(3,2)--(3,0)--(5.5,0);\node[below]at(3,0){$\varepsilon_F$};\node[above]at(3,2.8){$T=0$ 的填充基态};
}{定粒子数费米气体的零温参考态；满占据以下粒子不能都独立获得经典热激发。}{advA2:quality:fermi-c-q}}
\newcommand{\AAFermiCapacity}{\ThermalFigure{
 \draw[->](0,0)--(6,0)node[right]{$\varepsilon$};\draw[->](0,0)--(0,2.5)node[above]{$f$};\draw[dashed](0,2)--(3,2)--(3,0)--(5.5,0);
 \draw[DeepBlue,thick,domain=0:5.7,samples=100]plot(\x,{2/(exp(3*(\x-3))+1)});\draw[<->](2.3,2.7)--node[above]{$O(kT)$}(3.7,2.7);\node[below]at(3,0){$\mu\simeq\varepsilon_F$};
}{只有费米面附近的窄能窗参与低温热激发，窗外分布几乎不变。}{advA2:quality:fermi-c-a}
图~\ref{advA2:quality:fermi-c-a}中的平滑阶跃还须保持总粒子数不变。索末菲展开分别给出
\[
 0=g(\varepsilon_F)\delta\mu+\frac{\pi^2}{6}(kT)^2g'(\varepsilon_F)+O(T^4),
\]
再对能量积分展开并消去 $\delta\mu$，只剩 $\delta U=(\pi^2/6)g(\varepsilon_F)(kT)^2$。这样明确了定 $N$ 条件的作用；若固定化学势，不能不加检查地使用同一展开式。
}
\newcommand{\AQBoseCapacity}{\ThermalFigure{
 \draw[thick](0,0)--(5,0);\draw[thick](0,1)--(5,1) (0,1.8)--(5,1.8) (0,2.4)--(5,2.4);
 \foreach\x in {.5,1,1.5,2,2.5}{\fill[DeepBlue](\x,0)circle(2pt);}\foreach\x/\y in {3/1,4/1.8,2.8/2.4}{\fill(\x,\y)circle(2pt);}
 \node[right]at(6,0){基态 $N_0$};\node[right]at(6,1.8){激发态 $N-N_0$};
}{三维均匀理想玻色气体，将基态占据与激发态连续谱分开。}{advA2:quality:bose-c-q}}
\newcommand{\AABoseCapacity}{\ThermalFigure{
 \draw[->](0,0)--(5.2,0)node[right]{$T/T_c$};\draw[->](0,0)--(0,2.5)node[above]{$C_V/(Nk)$};
 \draw[DeepBlue,thick,domain=.001:1,samples=80]plot({1.5*\x},{1.92567168*\x^(1.5)});
 \draw[DeepBlue,thick]plot coordinates{(1.500000,1.92567168)(1.515000,1.91801858)(1.537500,1.90697113)(1.575000,1.88963169)(1.650000,1.85850730)(1.800000,1.80760611)(2.025000,1.75118594)(2.250000,1.71034420)(2.625000,1.66310220)(3.000000,1.63138079)(3.750000,1.59211964)(4.500000,1.56923444)};
 \draw[dashed](1.5,0)--(1.5,2.1);\node[below]at(1.5,0){$1$};\node[below]at(3,0){$2$};\node[below]at(4.5,0){$3$};\draw[dashed](0,1.5)--(5,1.5)node[right]{$3/2$};
}{Bose积分给出的定粒子数热容；上支数值解粒子数方程，未用经验指数替代。}{advA2:quality:bose-c-a}
图~\ref{advA2:quality:bose-c-a}中的临界以上曲线须同时考虑逸度随温度改变。定义 $g_s(z)=\operatorname{Li}_s(z)$，在 $T>T_c$ 有
\[
 \frac NV=\lambda_T^{-3}g_{3/2}(z),\qquad
 \frac{C_V}{Nk}=\frac{15}{4}\frac{g_{5/2}(z)}{g_{3/2}(z)}
 -\frac94\frac{g_{3/2}(z)}{g_{1/2}(z)}.
\]
定 $N,V$ 求导时，利用 $z\,g_s'(z)=g_{s-1}(z)$ 消去 $\dd z/\dd T$，便得到第二项。$z\to1^-$ 时 $g_{1/2}$ 发散，第二项趋零，与下支连续相接；$z\to0$ 时各 $g_s\simeq z$，回到经典极限 $3/2$。
}
\newcommand{\AQQuantumVirial}{\ThermalFigure{
 \foreach\x in {0,3,6,9}{\fill[DeepBlue](\x,0)circle(2pt);\draw[dashed](\x,0)circle(.7);}
 \draw[<->](0,-1.3)--node[below]{平均间距 $n^{-1/3}$}(3,-1.3);\draw[<->](-.7,1.2)--node[above]{$\lambda_T$}(.7,1.2);
}{热波长与平均间距决定占据稀疏程度，内禀简并度 $g$ 还会分摊每态粒子数。}{advA2:quality:virial-q}}
\newcommand{\AAQuantumVirial}{\ThermalFigure{
 \draw[->](0,0)--(5,0)node[right]{$x=n\lambda_T^3/g$};\draw[->](0,0)--(0,3.5)node[above]{$p/(nkT)$};\node[left]at(0,1.8){$1$};
 \draw[dashed](0,1.8)--(4.5,1.8);\draw[DeepBlue,thick](0,1.8)--(4.5,2.8)node[right]{费米};\draw[thick](0,1.8)--(4.5,.8)node[right]{玻色};
 \node[above]at(3,3.6){只示意 $x\ll1$ 的一阶斜率，纵向修正已放大};
}{交换统计使弱简并压强分别高于或低于经典值，但不是粒子间真实势力。}{advA2:quality:virial-a}
图~\ref{advA2:quality:virial-a}中的两直线只在原点附近表示第二维里修正。将逸度展开反解成密度展开后再比较 $p$，才能确定修正系数；若直接把 $z$ 当作 $n\lambda_T^3/g$，会漏掉同阶的反解项，并把正确的 $2^{-5/2}$ 系数算错。
}

\section{统计力学}

\subsection{相空间、系综与配分函数}

考虑不显含时间的拉格朗日系统 $L(q_\alpha,\dot q_\alpha)$。定义
\[
 p_\alpha=\frac{\partial L}{\partial\dot q_\alpha},
 \qquad
 H(q,p)=\sum_\alpha p_\alpha\dot q_\alpha-L,
\]
哈密顿方程为
\[
 \dot q_\alpha=\frac{\partial H}{\partial p_\alpha},
 \qquad
 \dot p_\alpha=-\frac{\partial H}{\partial q_\alpha}.
\]
把 $f$ 个广义坐标和 $f$ 个广义动量合成相坐标
$\vect\xi=(q_1,\ldots,q_f,p_1,\ldots,p_f)$，则
\[
 \dot{\vect\xi}=\mathsf J\nabla_\xi H,
 \qquad
 \mathsf J=\begin{pmatrix}0&\mathsf1\\-\mathsf1&0\end{pmatrix}.
\]
与量子态计数协调的无量纲相体积元取为
$\dd\Omega=\dd^f q\,\dd^f p/h^f$。系综是在相空间内沿哈密顿轨道运动的大量代表点，其归一化密度记为 $\rho(\vect\xi,t)$。

\begin{advtheorem}{刘维尔定理}\label{advA2:thm:liouville}
哈密顿相流不可压缩，沿任一相轨道
\[
 \boxed{\frac{D\rho}{Dt}=0}.
\]
等价地，相空间任意随流体积保持不变。
\end{advtheorem}
证明：相空间概率守恒给出连续性方程
\[
 \frac{\partial\rho}{\partial t}
 +\nabla_\xi\cdot(\rho\dot{\vect\xi})=0.
\]
哈密顿方程使相流散度为
\[
 \nabla_\xi\cdot\dot{\vect\xi}
 =\sum_\alpha\left(
 \frac{\partial^2H}{\partial q_\alpha\partial p_\alpha}
 -\frac{\partial^2H}{\partial p_\alpha\partial q_\alpha}
 \right)=0.
\]
因此连续性方程化为
$\partial_t\rho+\dot{\vect\xi}\cdot\nabla_\xi\rho=D\rho/Dt=0$。这也补足了源资料中省略的证明。

玻尔兹曼在分子混沌等动力学假设下得到粗粒化的 $H$ 定理，
$H_{\rm B}=\int\rho\ln\rho\,\dd\Omega$ 随时间不增。要注意：刘维尔定理保证孤立哈密顿体系的\emph{细粒化}吉布斯熵
$-k\int\rho\ln\rho\,\dd\Omega$ 不变；不可逆熵增来自粗粒化或约化描述，二者并不矛盾。

微正则系综固定 $N,E,V$。令
\[
 \omega(N,E,V)=\int_{E\le H<E+\Delta E}\dd\Omega,
 \qquad
 \rho=\begin{cases}1/\omega,&E\le H<E+\Delta E,\\0,&\text{其余},\end{cases}
\]
并记 $\Omega_E=\int_{H\le E}\dd\Omega$，则窄能壳中
$\omega\simeq(\partial\Omega_E/\partial E)_{N,V}\Delta E$。系综平均为
$\langle A\rangle=\int A\rho\,\dd\Omega$，熵为
$S=-k\int\rho\ln\rho\,\dd\Omega=k\ln\omega$。

\begin{advexamplebox}{等概率分布使微正则熵最大}\label{advA2:ex:microcanonical-maximum}
证明：在给定能壳和归一化条件下，均匀分布 $\rho=1/\omega$ 使
$S=-k\int\rho\ln\rho\,\dd\Omega$ 最大。

\AQMicrocanonical
\end{advexamplebox}
\begin{advsolution}
\AAMicrocanonical

对任意归一化分布 $\rho$，与均匀分布 $\rho_0=1/\omega$ 的相对熵满足
\[
 D(\rho\Vert\rho_0)=\int\rho\ln\frac{\rho}{\rho_0}\dd\Omega\ge0.
\]
于是
\[
 -\int\rho\ln\rho\,\dd\Omega
 \le-\int\rho\ln\rho_0\,\dd\Omega=\ln\omega.
\]
等号只在 $\rho=\rho_0$ 几乎处处成立时取得。也可用拉格朗日乘子变分，二阶变分
$-k\int(\delta\rho)^2/(2\rho)\dd\Omega<0$ 保证是最大值。
\end{advsolution}

两个弱耦合体系交换能量时，总态数
$\omega_1(E_1)\omega_2(E-E_1)$ 在平衡处最大，给出
\[
 \frac1{T_i}=\left(\frac{\partial S_i}{\partial E_i}\right)_{V_i,N_i},
 \qquad T_1=T_2.
\]
同理，允许体积和粒子交换时分别得到
$(\partial S/\partial V)_{E,N}=p/T$、
$(\partial S/\partial N)_{E,V}=-\mu/T$，以及
$p_1=p_2,\mu_1=\mu_2$，从微观态计数恢复
$\dd E=T\dd S-p\dd V+\mu\dd N$。

小体系与大热库接触时，把热库熵在总能量附近展开：
\[
 \ln\omega_B(E_T-E)=\text{常数}-\frac{E}{kT}+\cdots,
\]
故正则系综权重为 $e^{-E/(kT)}$。若还与粒子库接触，展开中出现化学势项，巨正则权重为
$e^{-(E-\mu N)/(kT)}$。

\begin{advdefinition}{配分函数与巨配分函数}\label{advA2:def:partition-functions}
\[
 Z(N,V,T)=\int e^{-\beta H_N}\dd\Omega_N,
 \qquad
 \Xi(\mu,V,T)=\sum_{N=0}^\infty e^{\beta\mu N}Z(N,V,T),
 \quad \beta=\frac1{kT}.
\]
热力学势为
\[
 F=-kT\ln Z,
 \qquad
 \Omega=-kT\ln\Xi=-pV.
\]
\end{advdefinition}
由此
\[
 U=-\frac{\partial\ln Z}{\partial\beta},
 \quad
 \mu=F_N,\quad p=-F_V,\quad S=-F_T.
\]

\begin{advexamplebox}{由巨配分函数直接求热力学量}\label{advA2:ex:grand-observables}
用 $\Xi(\mu,V,T)$ 写出平均粒子数、压强、熵和内能。

\AQGrand
\end{advexamplebox}
\begin{advsolution}
\AAGrand

对 $z=e^{\beta\mu}$ 求导，
\[
 \boxed{\langle N\rangle
 =z\frac{\partial\ln\Xi}{\partial z}
 =kT\left(\frac{\partial\ln\Xi}{\partial\mu}\right)_{T,V}}.
\]
由 $\Omega=-kT\ln\Xi$ 及 $p=-(\partial\Omega/\partial V)_{T,\mu}$，
\[
 \boxed{p=kT\left(\frac{\partial\ln\Xi}{\partial V}\right)_{T,\mu}},
 \qquad
 \boxed{S=k\ln\Xi+kT
 \left(\frac{\partial\ln\Xi}{\partial T}\right)_{\mu,V}}.
\]
固定 $z,V$ 对 $\beta$ 求导得到
\[
 \boxed{U=-\left(\frac{\partial\ln\Xi}{\partial\beta}\right)_{z,V}},
\]
等价形式是在固定 $\mu$ 时
$U-\mu N=-(\partial\ln\Xi/\partial\beta)_{\mu,V}$。最后
$U=TS-pV+\mu N$ 可作一致性校验。
\end{advsolution}

\subsection{经典统计结论}

无内禀简并的自由粒子单粒子配分函数写成
\[
 z_1=\int_0^\infty g(\varepsilon)e^{-\beta\varepsilon}\dd\varepsilon.
\]

\begin{advexamplebox}{不同维数自由粒子的态密度}\label{advA2:ex:density-of-states}
求一、二、三维非相对论自由粒子和无质量相对论粒子的态密度 $g(\varepsilon)$。

\AQDensityStates
\end{advexamplebox}
\begin{advsolution}
\AADensityStates

每个平动量子态占据相空间体积 $(2\pi\hbar)^d$。先数动量球内态数，再对能量求导。
对 $\varepsilon=p^2/(2m)$，
\[
 \boxed{g_{3}(\varepsilon)=\frac{V(2m)^{3/2}}{4\pi^2\hbar^3}\varepsilon^{1/2}},
\]
\[
 \boxed{g_{2}(\varepsilon)=\frac{Am}{2\pi\hbar^2}},
 \qquad
 \boxed{g_{1}(\varepsilon)=\frac{L}{\pi\hbar}\sqrt{\frac{m}{2\varepsilon}}}.
\]
对 $\varepsilon=pc$，
\[
 \boxed{g_{3}^{\rm rel}=\frac{V\varepsilon^2}{2\pi^2\hbar^3c^3}},
 \quad
 \boxed{g_{2}^{\rm rel}=\frac{A\varepsilon}{2\pi\hbar^2c^2}},
 \quad
 \boxed{g_{1}^{\rm rel}=\frac{L}{\pi\hbar c}}.
\]
每一式的 $g\dd\varepsilon$ 都无量纲；若存在自旋简并度，还应整体乘以简并因子。
\end{advsolution}

\begin{advexamplebox}{自由粒子配分函数与物态方程}\label{advA2:ex:free-particle-partition}
计算三维非相对论和无质量相对论自由粒子的单粒子配分函数，并证明平均能量分别为
$3kT/2$、$3kT$，且分别满足 $pV=2E/3$、$pV=E/3$。

\AQFreePartition
\end{advexamplebox}
\begin{advsolution}
\AAFreePartition

令热德布罗意波长
\[
 \lambda_T=\sqrt{\frac{2\pi\hbar^2}{mkT}}.
\]
高斯积分给出非相对论单粒子
$z_1=V/\lambda_T^3\propto VT^{3/2}$；无质量粒子由
$\int_0^\infty\varepsilon^2e^{-\beta\varepsilon}\dd\varepsilon=2/\beta^3$ 得
\[
 z_1^{\rm rel}=\frac{V(kT)^3}{\pi^2\hbar^3c^3}\propto VT^3.
\]
于是 $E_1=-\partial_\beta\ln z_1$ 分别为 $3kT/2$ 和 $3kT$。两者都因 $\ln z_1$ 含 $\ln V$ 而给出
$p=kT(\partial\ln z_1/\partial V)_T=kT/V$，所以
\[
 \boxed{pV=\frac23E}\quad\text{和}\quad
 \boxed{pV=\frac13E}.
\]
这里 $E$ 可指单粒子平均能量；乘以粒子数后关系不变。
\end{advsolution}

经典全同粒子必须除去置换重复计数：$Z_N=z_1^N/N!$。斯特林公式
$\ln N!\simeq N\ln N-N$ 使自由能和熵恢复广延性。

\begin{advexamplebox}{单原子经典理想气体熵}\label{advA2:ex:classical-entropy}
由 $Z_N=(V/\lambda_T^3)^N/N!$ 证明
\[
 S=\frac32Nk\ln T+Nk\ln\frac VN+Nks_0,
 \quad
 s_0=\frac32\ln\frac{mk}{2\pi\hbar^2}+\frac52.
\]

\AQClassicalEntropy
\end{advexamplebox}
\begin{advsolution}
\AAClassicalEntropy

用 $F=-kT\ln Z_N$、$S=-(\partial F/\partial T)_{V,N}$：
\[
 S=Nk\left[\ln\frac{V}{N\lambda_T^3}+\frac52\right].
\]
展开 $\lambda_T$ 便得到题式。它与前面的萨克--泰特洛德结果完全一致；$1/N!$ 若遗漏，会多出非广延的 $Nk\ln N$，这就是吉布斯佯谬的经典表现。
\end{advsolution}

\begin{advexamplebox}{$N$ 个一维经典谐振子}\label{advA2:ex:harmonic-partition}
计算 $N$ 个频率为 $\omega$、彼此独立且固定在不同位置的一维经典谐振子的配分函数、内能与热容。

\AQOscillators
\end{advexamplebox}
\begin{advsolution}
\AAOscillators

单振子的哈密顿量 $H=p^2/(2m)+m\omega^2q^2/2$，故
\[
 z=\frac1h\int_{-\infty}^{\infty}\!\dd p\dd q\,e^{-\beta H}
 =\frac{kT}{\hbar\omega}.
\]
振子依附在不同位置，彼此可分辨，不除以 $N!$：
\[
 \boxed{Z=z^N=\left(\frac{kT}{\hbar\omega}\right)^N},
 \quad
 \boxed{U=NkT},\quad \boxed{C_V=Nk}.
\]
两个二次项各贡献 $kT/2$，与能量均分定理一致。低温时经典结果失效，量子配分函数应取
$e^{-\beta\hbar\omega/2}/(1-e^{-\beta\hbar\omega})$。
\end{advsolution}

\begin{advexamplebox}{$N$ 个三自由度刚体转子}\label{advA2:ex:rigid-rotor-partition}
主惯量为 $I_1,I_2,I_3$ 的非线性刚体分子只有转动自由度。求其高温经典配分函数，并给出转动内能与热容。

\AQRotors
\end{advexamplebox}
\begin{advsolution}
\AARotors

对欧拉角和三个共轭角动量作高斯积分，单分子转动配分函数为
\[
 z_{\rm rot}=\frac{\sqrt\pi}{\sigma}
 \left(\frac{8\pi^2kT}{h^2}\right)^{3/2}
 \sqrt{I_1I_2I_3},
\]
其中 $\sigma$ 是分子取向的对称数。若这些转子位于固定晶格点，$Z=z_{\rm rot}^N$；若作为全同自由分子的内部因子，则它进入
$Z=(z_{\rm trans}z_{\rm rot})^N/N!$。因 $z_{\rm rot}\propto T^{3/2}$，
\[
 \boxed{U_{\rm rot}=\frac32NkT},
 \qquad \boxed{C_{V,\rm rot}=\frac32Nk}.
\]
三项角动能各贡献 $kT/2$。
\end{advsolution}

\begin{advexamplebox}{匀强电场中的刚性偶极子}\label{advA2:ex:dipole-partition}
$N$ 个大小为 $d$ 的刚性电偶极子置于匀强电场 $E$ 中，忽略动能。求配分函数、平均偶极矩与弱场极化率。

\AQDipoleStatistics
\end{advexamplebox}
\begin{advsolution}
\AADipoleStatistics

取 $x=dE/(kT)$。单偶极子取向积分为
\[
 z_1=\int e^{x\cos\theta}\dd\Omega
 =4\pi\frac{\sinh x}{x},
 \qquad Z=z_1^N.
\]
沿场方向的平均偶极矩为
\[
 \langle d_\parallel\rangle
 =kT\frac{\partial\ln z_1}{\partial E}
 =d\mathcal L(x),
 \quad \mathcal L(x)=\coth x-\frac1x.
\]
故极化强度
$P=Nd\mathcal L(x)/V$。弱场 $\mathcal L(x)=x/3+O(x^3)$，
\[
 \boxed{\chi_e=\frac1{\varepsilon_0}\frac{\partial P}{\partial E}
 =\frac{Nd^2}{3\varepsilon_0VkT}}.
\]
强场 $x\to\infty$ 时 $P\to Nd/V$，正确再现饱和极化。
\end{advsolution}

\subsection{费米统计与玻色统计}

把单粒子态记为 $\varepsilon_i$、占据数记为 $n_i$，则
$E=\sum_i n_i\varepsilon_i,N=\sum_i n_i$。每个态的巨配分函数为
\[
 \Xi_i=\sum_{n_i}e^{-\beta(\varepsilon_i-\mu)n_i}.
\]
费米子有 $n_i=0,1$；玻色子有 $n_i=0,1,2,\ldots$。

\begin{advexamplebox}{费米--狄拉克与玻色--爱因斯坦分布}\label{advA2:ex:quantum-distributions}
求两种统计的平均占据数和巨势，并写出压强公式。

\AQQuantumDistributions
\end{advexamplebox}
\begin{advsolution}
\AAQuantumDistributions

令 $z_i=e^{-\beta(\varepsilon_i-\mu)}$。费米子
$\Xi_i=1+z_i$，玻色子 $\Xi_i=(1-z_i)^{-1}$。用
$\langle n_i\rangle=z_i\partial_{z_i}\ln\Xi_i$，
\[
 \boxed{\langle n_i\rangle=\frac1{e^{(\varepsilon_i-\mu)/(kT)}\pm1}},
\]
上号对应费米、下号对应玻色。又
\[
 \boxed{\Omega=\mp kT\sum_i\ln\!\left(1\pm e^{-\beta(\varepsilon_i-\mu)}\right)},
\]
因此
\[
 \boxed{pV=\pm kT\sum_i\ln\!\left(1\pm e^{-\beta(\varepsilon_i-\mu)}\right)}.
\]
有限系统的玻色巨配分级数严格要求 $\mu$ 低于基态能量（凝聚热力学极限另行处理）；费米式自动体现泡利不相容原理。
\end{advsolution}

三维理想费米气在零温下把单粒子态填充到费米能
\[
 \varepsilon_F=\frac{\hbar^2}{2m}
 \left(\frac{6\pi^2N}{gV}\right)^{2/3},
 \qquad p_F=\frac{2N\varepsilon_F}{5V},
\]
其中 $g$ 是内禀简并度。

\begin{advexamplebox}{金属电子的低温热容}\label{advA2:ex:fermi-heat-capacity}
证明 $T\ll T_F$ 时理想费米气的电子热容为
\[
 \boxed{C_V=\frac{\pi^2}{2}Nk\frac{kT}{\varepsilon_F}}.
\]

\AQFermiCapacity
\end{advexamplebox}
\begin{advsolution}
% CompletenessThermalFields20261024: fixed-N Sommerfeld derivation.
先说明低温展开怎样产生，再实施定粒子数约束。定义 $T_F=\varepsilon_F/k$，以下固定 $N,V$，
并令 $g(\varepsilon)$ 表示包含全部自旋简并的\emph{总态密度}，单位为能量的倒数。
费米分布为 $f_T(\varepsilon)=[\exp((\varepsilon-\mu)/(kT))+1]^{-1}$。
图~\ref{advA2:complete:sommerfeld}把它与同一化学势处的零温阶跃相减；这一步把整个费米海的积分
变成只在 $|\varepsilon-\mu|\sim kT$ 内有贡献的修正。

\begin{qualityfigure}
\centering
\begin{tikzpicture}[>=Latex,font=\small,x=1.0cm,y=2.8cm]
 \draw[->](-4.6,0)--(4.6,0)node[right]{$x=(\varepsilon-\mu)/(kT)$};
 \draw[->](0,-.62)--(0,.69)node[above]{$D(x)$};
 \fill[DeepBlue!13,domain=0:4,samples=70] (0,0)--plot(\x,{1/(exp(\x)+1)})--(4,0)--cycle;
 \fill[DeepBlue!13,domain=-4:0,samples=70] (-4,0)--plot(\x,{-1/(exp(-\x)+1)})--(0,0)--cycle;
 \draw[DeepBlue,thick,domain=0:4,samples=70]plot(\x,{1/(exp(\x)+1)});
 \draw[DeepBlue,thick,domain=-4:0,samples=70]plot(\x,{-1/(exp(-\x)+1)});
 \draw[dashed](0,-.5)--(0,.5);
 \node[above,align=center]at(2,.25){增占据\\$x>0$};
 \node[below,align=center]at(-2,-.25){减占据\\$x<0$};
 \node[below]at(3.7,-.47){$D(-x)=-D(x)$};
 \node[above]at(-2.6,.48){$\int D(x)\,\dd x=0$};
\end{tikzpicture}
\figcap{费米分布相对零温阶跃的差 $D(x)=(e^x+1)^{-1}-\Theta(-x)$。
两瓣的有向面积相消，但乘上 $x$ 后两侧贡献同号；这决定首项为 $T^2$ 修正。\label{advA2:complete:sommerfeld}}
\end{qualityfigure}

对在 $\mu$ 附近光滑、无带边或奇点的函数 $\Phi$，考虑
\[
 I_\Phi=\int_0^\infty\Phi(\varepsilon)f_T(\varepsilon)\,\dd\varepsilon.
\]
设 $\mu\gg kT$。作 $x=(\varepsilon-\mu)/(kT)$ 代换，并在费米面附近展开 $\Phi$，得到
\[
 I_\Phi-\int_0^\mu\Phi(\varepsilon)\,\dd\varepsilon
 =kT\int_{-\mu/(kT)}^\infty\Phi(\mu+kTx)D(x)\,\dd x.
\]
由于 $D$ 的尾部指数衰减，各项的矩积分可把下限延到 $-\infty$，所添误差指数小。
图中 $D$ 为奇函数，因此常数项和二次项不贡献；一次项的系数为
\[
 \int_{-\infty}^{\infty}xD(x)\,\dd x
 =2\int_0^\infty\frac{x\,\dd x}{e^x+1}
 =2\sum_{r=1}^\infty\frac{(-1)^{r+1}}{r^2}
 =\frac{\pi^2}{6}.
\]
这里先展开 $(e^x+1)^{-1}=\sum_{r\ge1}(-1)^{r+1}e^{-rx}$，逐项积分用
$\int_0^\infty xe^{-rx}\dd x=r^{-2}$；交错平方和是普通平方和 $\pi^2/6$ 的一半。
于是索末菲展开的前两项为
\begin{equation}
 I_\Phi=\int_0^\mu\Phi(\varepsilon)\,\dd\varepsilon
 +\frac{\pi^2}{6}(kT)^2\Phi'(\mu)+O(T^4).
 \label{advA2:eq:complete-sommerfeld}
\end{equation}
余项表示在费米面光滑、导数有界的低温渐近展开；不能用在 $kT$ 已与 $\varepsilon_F$ 同量级的区域。

\textbf{先固定粒子数。} 取 $\Phi=g$，写 $\mu=\varepsilon_F+\delta\mu$。由于
$N=\int_0^{\varepsilon_F}g\,\dd\varepsilon$，式~\eqref{advA2:eq:complete-sommerfeld}给出
\[
 0=g_F\delta\mu+\frac{\pi^2}{6}(kT)^2g'_F+O(T^4),
 \qquad
 \delta\mu=-\frac{\pi^2}{6}(kT)^2\frac{g'_F}{g_F}+O(T^4),
\]
下标 $F$ 表示在 $\varepsilon_F$ 处取值。三维自由粒子有 $g=A\sqrt\varepsilon$，因此
\[
 g_F=\frac{3N}{2\varepsilon_F},\qquad
 \mu=\varepsilon_F-\frac{\pi^2}{12}\frac{(kT)^2}{\varepsilon_F}+O(T^4).
\]
化学势的微小下移正是补偿图中态密度不恒定时两侧粒子数不完全相等的修正。

\textbf{再计算能量。} 取 $\Phi=\varepsilon g(\varepsilon)$，保留同阶的两项：
\begin{align*}
 U(T)-U(0)
 &=\varepsilon_Fg_F\delta\mu
   +\frac{\pi^2}{6}(kT)^2(g_F+\varepsilon_Fg'_F)+O(T^4)\\
 &=-\frac{\pi^2}{6}(kT)^2\varepsilon_Fg'_F
   +\frac{\pi^2}{6}(kT)^2(g_F+\varepsilon_Fg'_F)+O(T^4)\\
 &=\frac{\pi^2}{6}g_F(kT)^2+O(T^4)
  =\frac{\pi^2}{4}N\frac{(kT)^2}{\varepsilon_F}+O(T^4).
\end{align*}
第一项来自积分上限的移动；若误把 $\mu$ 固定，就会漏掉这一项并改变系数。
最后在固定 $N,V$ 下求导，得到
\[
 \boxed{C_V=\frac{\pi^2}{2}Nk\frac{kT}{\varepsilon_F}+O(T^3)}.
\]
$C_V$ 的单位为能量除以温度，$T\to0$ 时线性趋零；只有费米面附近的窄能窗能被热激发，
所以它比经典值 $3Nk/2$ 小一个 $T/T_F$ 因子。

\end{advsolution}

理想玻色气的激发态粒子数为
\[
 N_e=\int_0^\infty\frac{g(\varepsilon)\dd\varepsilon}
 {e^{(\varepsilon-\mu)/(kT)}-1}.
\]
三维均匀气体在临界温度以下有 $\mu=0$、
\[
 N_0=N\left[1-\left(\frac T{T_c}\right)^{3/2}\right].
\]

\begin{advexamplebox}{玻色--爱因斯坦凝聚的热容曲线}\label{advA2:ex:bec-heat-capacity}
定性并定量说明三维理想玻色气定容热容随温度的变化，并画出临界点附近曲线。

\AQBoseCapacity
\end{advexamplebox}
\begin{advsolution}
\AABoseCapacity

$T<T_c$ 时
\[
 U=\frac32kT\frac{V}{\lambda_T^3}\zeta(5/2),
\]
故
\[
 \boxed{C_V=\frac{15}{4}Nk
 \frac{\zeta(5/2)}{\zeta(3/2)}
 \left(\frac T{T_c}\right)^{3/2}}.
\]
它从零按 $T^{3/2}$ 增加，在 $T_c^-$ 达到约 $1.93Nk$。$T>T_c$ 时化学势转为负值，热容从同一数值连续接出，随后趋向经典极限 $3Nk/2$；一阶导数在 $T_c$ 处不连续，形成尖点而非潜热跃变。


图~\ref{advA2:quality:bose-c-a}的两支在临界点相接；基态在本模型中不贡献激发能，低温热容来自激发态占据随温度改变。
\end{advsolution}

\begin{advexamplebox}{稀薄量子气体的第二维里修正}\label{advA2:ex:quantum-virial}
对三维非相对论理想费米气和玻色气：
\begin{enumerate}
 \item 写出经典极限的化学势、内能和压强。
 \item 给出弱简并条件，并求到 $n\lambda_T^3/g$ 一阶的 $\mu,U,p$。
 \item 把物态方程与无排斥体积的范德瓦耳斯形式
 $(p+aN^2/V^2)V=NkT$ 比较，说明这种等价的物理边界。
\end{enumerate}

\AQQuantumVirial
\end{advexamplebox}
\begin{advsolution}
\AAQuantumVirial

令 $\eta=+1$ 表示费米统计、$\eta=-1$ 表示玻色统计，逸度
$z=e^{\mu/(kT)}$，并记 $x=n\lambda_T^3/g$。经典麦克斯韦--玻尔兹曼极限为
\[
 \boxed{\mu_{\rm MB}=kT\ln x},\qquad
 \boxed{U_{\rm MB}=\frac32NkT},\qquad
 \boxed{p_{\rm MB}=nkT}.
\]
弱简并条件是
\[
 \boxed{x=\frac{n\lambda_T^3}{g}\ll1},
\]
即平均每个热波包和每个内禀态中的粒子数远小于1。展开分布函数，
\[
 x=z-\eta\frac{z^2}{2^{3/2}}+O(z^3),
 \qquad
 \frac{p\lambda_T^3}{gkT}
 =z-\eta\frac{z^2}{2^{5/2}}+O(z^3).
\]
反解并保留一阶修正：
\[
 \begin{aligned}
 \mu&=kT\left[\ln x+\eta\frac{x}{2^{3/2}}\right],\\
 p&=nkT\left[1+\eta\frac{x}{2^{5/2}}\right],\qquad U=\frac32pV.
 \end{aligned}
\]
与 $p=nkT-an^2$ 比较，得到 $a(T)=-\eta kT\lambda_T^3/(2^{5/2}g)$。
费米统计给出正压修正，可看成“统计排斥”；玻色统计给出负压修正，可看成“统计聚集”。它只是第二维里系数的数学等价，并不意味着无相互作用量子粒子之间真的出现了范德瓦耳斯力；密度升高后高阶项不可忽略，等价立即失效。
\end{advsolution}


\clearpage\endgroup
